\section{Dataset thinking: если потом делать LoRA/DPO, какие примеры собирать}

Если в будущем потребуется дообучение (LoRA/DPO) поверх базовой модели, а не только prompt-инжиниринг, принципы сбора примеров те же, что и для few-shot exemplars (13.8), но с дополнительным требованием: разнообразие покрытия по архетипическому вектору (13.2), а не просто общее количество примеров.

Практически это значит: датасет должен содержать примеры для каждого режима из voice tests (13.5), для каждого stress test (13.4), и — что легко упустить — примеры отказа модели скатиться в failure modes главы 9, а не только позитивные примеры желаемого поведения. DPO-подход в особенности выигрывает от пар "хороший/плохой" ответ на одну и ту же ситуацию, где плохой ответ — конкретный, названный failure mode (9.1--9.8), а не произвольно испорченная версия хорошего.

Это не техническая спецификация обучения (она вне области этой книги), а требование к содержанию данных: датасет, собранный без привязки к архетипическому вектору, рискует научить модель звучать похоже на Рин, не научив её быть Рин в структурном смысле, установленном главами 2--12.
